\documentclass[11pt]{article}

\usepackage[margin=1in]{geometry}
\usepackage{amsmath,amssymb,mathtools}
\usepackage{booktabs,array,enumitem}
\usepackage[colorlinks=true,linkcolor=blue,urlcolor=blue,citecolor=blue]{hyperref}

\newcounter{algorithm}
\setlist{itemsep=0.35em,topsep=0.35em}

\newcommand{\R}{\mathbb{R}}
\newcommand{\norm}[1]{\left\lVert #1\right\rVert}
\newcommand{\grad}{\nabla}
\newcommand{\hess}{\nabla^2}
\newcommand{\lammin}{\lambda_{\min}}

\title{Implementation Details for the Paper's First-Order Adaptive UTR Method}
\author{Adapted from the adaptive UTR algorithm in arXiv:2311.11489}
\date{}

\begin{document}
\maketitle

\begin{quote}\small
This note keeps only the information needed to implement the adaptive
universal trust-region method described in the paper when the goal is an
approximately first-order stationary point.  The second-order phase,
assumptions, lemmas, complexity bounds, and convergence proofs are
intentionally omitted.
\end{quote}

\vspace{0.5em}

\section{Paper implementation (only for getting small gradient)}

\subsection{Quantities used by the method}

At outer iteration $k$, define
\[
    g_k=\grad f(x_k),
    \qquad
    H_k=\hess f(x_k).
\]
Given a tolerance $\epsilon>0$, stop and return $x_k$ if
\[
    \norm{g_k}<\epsilon.
\]
Thus every trial below is computed only while $\norm{g_k}\geq\epsilon$.  The
Hessian is still used to select the regularization and trust-region parameters,
but it is not part of the stopping test.

The method maintains a persistent penalty parameter $\rho_k>0$.  At the start
of the inner loop, set
\[
    \rho_k^{(0)}=\rho_k.
\]
The inner-loop value $\rho_k^{(j)}$ is used to choose the regularization
parameter $\sigma_k^{(j)}$ and radius scale $r_k^{(j)}$.

\subsection{Acceptance tests}

The method accepts an inner-loop trial when the monotonicity condition
\[
    f(x_k+d_k^{(j)})-f(x_k)\leq 0
\]
holds and at least one of the following first-order conditions holds:
\[
    f(x_k+d_k^{(j)})-f(x_k)
    \leq
    -\frac{\eta}{\rho_k^{(j)}}\norm{g_k}^{3/2}
    \quad\text{or}\quad
    \norm{\grad f(x_k+d_k^{(j)})}\leq \xi\norm{g_k}.
\]

\subsection{Parameter-selection strategy}

The following strategy chooses $(\sigma_k^{(j)},r_k^{(j)})$ from the current
inner penalty value $\rho_k^{(j)}$ while $\norm{g_k}\geq\epsilon$.

\begin{enumerate}[leftmargin=2.2em,label=\textbf{Case \arabic*.}]
    \item If
    \[
        \left|\lammin(H_k)\right|
        \geq \rho_k^{(j)}\norm{g_k}^{1/2},
    \]
    choose
    \[
        \sigma_k^{(j)}=0,\qquad
        r_k^{(j)}=\frac{1}{2\rho_k^{(j)}}.
    \]

    \item If
    \[
        \left|\lammin(H_k)\right|
        < \rho_k^{(j)}\norm{g_k}^{1/2},
    \]
    choose
    \[
        \sigma_k^{(j)}=\rho_k^{(j)},\qquad
        r_k^{(j)}=\frac{1}{4\rho_k^{(j)}}.
    \]
\end{enumerate}

\paragraph{Remark (avoiding an explicit eigenvalue calculation).}
Let
\[
    \tau_k^{(j)}=\rho_k^{(j)}\norm{g_k}^{1/2}.
\]
Because $H_k$ is symmetric, the two shifted matrices satisfy
\[
    H_k-\tau_k^{(j)}I\succ 0
    \iff \lammin(H_k)>\tau_k^{(j)},
    \qquad
    H_k+\tau_k^{(j)}I\succ 0
    \iff \lammin(H_k)>-\tau_k^{(j)}.
\]
Consequently, the minimum eigenvalue need not be computed explicitly.  Attempt
Cholesky factorizations of the shifted matrices.  Use Case~1 if
$H_k+\tau_k^{(j)}I$ is not positive definite or if
$H_k-\tau_k^{(j)}I$ is positive definite; otherwise use Case~2.  A failed
Cholesky factorization detects a nonpositive pivot.  This test is often cheaper
than an eigendecomposition and can exploit sparse factorizations.

\subsection{Algorithmic form}

\begin{center}
\begin{minipage}{0.96\linewidth}
\hrule
\vspace{0.4em}
\refstepcounter{algorithm}\label{alg:paper-adaptive-utr}
\noindent\textbf{Algorithm~\thealgorithm: First-order adaptive universal trust-region method}

\smallskip
\noindent\textbf{Data:} Initial point $x_0$, tolerance $\epsilon>0$,
acceptance constants $0<\eta<1/32$ and $1/4<\xi<1$, initial penalty
$\rho_0>0$, minimum penalty $\rho_{\min}>0$, penalty increase factor
$\mu_1>1$, penalty decrease factor $\mu_2>1$, and maximum number of inner
trials per outer iteration $J_{\max}\in\mathbb{N}$.

\begin{enumerate}[leftmargin=2.2em,label=\arabic*.]
    \item For $k=0,1,\ldots$:
    \begin{enumerate}[leftmargin=2.4em,label*=\arabic*.]
        \item Compute $g_k=\grad f(x_k)$.  If
        $\norm{g_k}<\epsilon$, output $x_k$ and terminate.  Otherwise, compute
        $H_k=\hess f(x_k)$ and set $\rho_k^{(0)}=\rho_k$.

        \item For $j=0,\ldots,J_{\max}-1$:
        \begin{enumerate}[leftmargin=2.4em,label*=\arabic*.]
            \item Choose $(\sigma_k^{(j)},r_k^{(j)})$ using the
            parameter-selection strategy above.

            \item Solve the shared trust-region subproblem
            using $(\sigma_k,r_k)=(\sigma_k^{(j)},r_k^{(j)})$ and obtain
            $d_k^{(j)}$.

            \item If the monotonicity condition and the first-order acceptance
            test hold, accept the inner-loop trial and break.

            \item Otherwise reject the trial and set
            \[
                \rho_k^{(j+1)}=\mu_1\rho_k^{(j)}.
            \]
        \end{enumerate}

        \item If no trial was accepted, output $x_k$ and terminate with status
        \texttt{INNER\_ITERATION\_LIMIT}. Otherwise, update
        \[
            x_{k+1}=x_k+d_k^{(j)},
            \qquad
            \rho_{k+1}
            =
            \max\left\{\rho_{\min},\frac{\rho_k^{(j)}}{\mu_2}\right\}.
        \]
    \end{enumerate}
\end{enumerate}
\vspace{0.1em}
\hrule
\end{minipage}
\end{center}

\subsection{Default algorithmic parameters}

$\rho_0 = 1$, $\eta = 29/1200$, $\rho_{\min} = 10^{-5}$, $\xi = 0.5$, $\mu_1 = 2.25$, $\mu_2 = 1.75$, $\epsilon = 10^{-5}$, $J_{\max}=100$

\end{document}
